\fbox{
\parbox{0.48\textwidth}{
\rule{0.48\textwidth}{0.4pt}
{\bf Context}: There is a major force uniting America’s fiercely partisan politicians: big technology companies. Democrats and Republicans at the federal and state levels are coming together to scrutinize the power of the Silicon Valley giants and, potentially, to rein them in. Letitia James, the Democratic attorney general of New York, announced on Friday that attorneys general in eight states — four Democrats and four Republicans — and the District of Columbia had begun an antitrust investigation of Facebook.

\rule{0.48\textwidth}{0.4pt}
{\bf Megatron-LM}: The probe, which Ms. James said was likely to result in a lawsuit, is the latest in a string of high-profile regulatory actions, in state and federal courts, that challenge the market dominance of tech companies. The Justice Department is also investigating Facebook's data practices and whether the company violated a 2011 consent decree with the Federal Trade Commission that required the company to obtain explicit consent from users before sharing their personal information with third parties.
\\[1\baselineskip]
What might prompt a bipartisan coalition to sue Facebook? It isn't clear that there's any one thing. But lawmakers and regulators have been increasingly alarmed by the company's recent controversies, which include a data breach affecting more than 50 million users, the spread of fake news on the platform, and the spreading of propaganda by Russian-backed accounts during the 2016 election.
\\[1\baselineskip]
It's also not clear how serious the Democratic attorneys general are about filing a lawsuit. While the attorneys general announced their investigation in a press release, they said they had not yet reached a decision on whether to sue. In the past, some of the attorneys general have focused their efforts on trying to extract settlements from the tech companies without filing lawsuits.
\\[1\baselineskip]
"As attorneys general who have investigated and prosecuted health-care fraud, price-fixing, and other violations of our states' antitrust laws, we are committed to holding companies accountable when they fail to protect consumers and competition in the market," Ms. James said in a statement.
\\[1\baselineskip]
The investigation, which was first reported by The New York Times, was announced as the companies that dominate the tech industry are facing scrutiny from antitrust regulators around the world. Google, Twitter, Facebook and Amazon have faced a series of antitrust probes and class-action lawsuits in the United States and Europe ... 
%Last month, European regulators fined Google \$2.7 billion for unfairly promoting its own shopping service. And the European Union's competition commissioner, Margrethe Vestager, has threatened to hit Google with additional penalties.

}}

\fbox{%
\parbox{0.48\textwidth}{%
\rule{0.48\textwidth}{0.4pt}
{\bf Context}: A flower, sometimes known as a bloom or blossom, is the reproductive structure found in flowering plants (plants of the division Magnoliophyta, also called angiosperms). The biological function of a flower is to effect reproduction, usually by providing a mechanism for the union of sperm with eggs.

\rule{0.48\textwidth}{0.4pt}
{\bf Megatron-LM}: The structural components of the flower are the petals and sepals, which are the reproductive organs, the hypanthium, which forms the base of the flower, and the calyx, which is the stigma. A flower is a composite structure, made up of sepals and petals, the sepals being composed of three fused tepals and the petals of two fused petals, though there are a few exceptions to this rule. The three sepals are fused to form a tube called the calyx, which is the structure that attracts the pollinating insects. The petals are composed of two fused sepals (the outer ones) which open and fold inward (clasp) at the tips, forming the petals. The four sepals and petals are united to form a corolla, the outermost structure of a flower that carries pollen.
\\[1\baselineskip]
The structure of the flower has important functional significance, including protection of the ovary from insects and the transfer of pollen between the male and female parts of a flower. The structure of the flower is also used by the plants to attract pollinators, for example by the shape of the petals and the colours of the stamens.
\\[1\baselineskip]
The structure of a flower varies widely between plants. The simplest flowers are unbranched and made up of only two sepals and a single petal. Other examples are composed of several petals joined together, with the innermost sepals completely fused to form a tube called the calyx. The petals are usually coloured to attract insects, which can only see certain wavelengths of light. The sepals are often coloured too, but usually have a more limited range of colours.
\\[1\baselineskip]
Some flowers have additional structures, such as stamens, anthers and ovaries. The presence of stamens and anthers in some flowers is used for cross-pollination. The ovary is the organ where the female gametes are produced, and is the first structure that is developed in an embryo of a flower. The ovary is the female flower's reproductive system, and consists of the ovary wall, the ovule, and the oviduct.
\\[1\baselineskip]
Flowers vary in their size and shape, from tiny to huge, and are often brightly coloured. Flowers may be seen as part of a larger structure, such as a flower head, an inflorescence, a bract, a spathe, or a carpels.
\\[1\baselineskip]
The term "flower" is usually applied to the floral structures that are visible externally, but many ...
%flowers produce their flowers internally, and these structures are collectively known as the "floral anatomy" or "floral anatomy of a plant". The floral anatomy includes structures that are only visible internally, such as the ovary, the style, and the stigma.
%\newline \newline
%The term flower is also applied to the fruit of a flowering plant, such as the berry of a strawberry. The term may also be used to refer to the various parts
%of the flower that are not visible externally.
%The term "flower" is also used in other contexts, including the common names of certain plants and the eponymous genus of an entire family of flowering plants, Asteraceae.

}}

\fbox{%
\parbox{0.48\textwidth}{%
\rule{0.48\textwidth}{0.4pt}
{\bf Context}: OpenAI did not release the full GPT-2 model due to concerns of malicious use, but they did release a smaller version equivalent in size to the original GPT (117 M parameters), trained on the new, larger dataset. Although not as powerful as the large model, the smaller version still has some language generation chops.

\rule{0.48\textwidth}{0.4pt}
{\bf Megatron-LM}: In December 2016, DeepMind released an even larger model, GPT-3, with more than 1.5 billion parameters. DeepMind has no plans to release the model, so the question remains whether other AI researchers will be able to copy and build on it.
\\[1\baselineskip]
Other language models have been built using other techniques. In 2016, a team of researchers at Columbia University, the University of Toronto, and Google DeepMind showed that neural network-based models, including deep neural networks, can learn the grammatical rules of a language. A group at Google DeepMind demonstrated that these rules could be represented in a network as a kind of decision tree, where the nodes in the tree are the words in the sentence, and the rules for the words are represented by the decision nodes at the end of the tree. The network would go through the network in order to choose the next word in the sentence based on whether it follows the rules or not. In 2017, the team at DeepMind released a paper detailing the training and performance of their model.
\\[1\baselineskip]
The problems with the model are many. One of the biggest issues is that grammatical rules are ambiguous. For example, the phrase "I came down" could mean I came down from the stairs, or I came down from the sky. In English, there are numerous other examples of sentences that are ambiguous. The model does not have any way to tell the difference between the ambiguous sentences, and so it simply makes a prediction about the entire sentence based on what rules it does know. This makes the model vulnerable to the phenomenon of overfitting. In order to avoid this, models must be trained using a large set of data, so that the model can find useful patterns in the data, and use those patterns to help make predictions. This approach is also prone to errors, as it relies on the training data to help make correct predictions, and the data might be incorrectly labeled, resulting in a model that produces a poor set of predictions. In the case of GPT-2, the model only has about 100 training sentences for each type of phrase (subject-verb-object).
\\[1\baselineskip]
Additionally, languages have different rules for different types of phrasing. In English, the subject-verb-object construction is fairly universal, but other languages have different rules, which makes training a GPT more complicated.
\\[1\baselineskip]
Finally, there is an inherent trade-off between the number of rules and the performance of the model ... 
%As a rule, the more rules the model has, the better it performs. However, as the number of rules increases, the likelihood of the model to overfit the training data increases, and the ability to perform generalizations decreases. For instance, if the model has 1,000 rules ... 
%for an English phrase, it will be able to predict other English phrases with high accuracy. However, it will struggle with generalizing to other languages. If the model has 10,000 rules, it will be able to perform generalizations, but it will be inaccurate with respect to English.
%There have been other grammar learning algorithms that have not yet been released to the public. One such system was used to produce a language translation system that was able to translate into and out of English. However, the system did not learn the rules for the different languages, but instead relied on information about the structure of the language to translate between the two languages. Another system has also been developed by Google DeepMind, which used neural networks to learn grammar. The system was able to learn grammatical rules for two languages, English and French.

}}


\fbox{%
\parbox{0.48\textwidth}{%
\rule{0.48\textwidth}{0.4pt}
{\bf Context}: The 36th International Conference on Machine Learning (ICML 2019) will be held in Long Beach, CA, USA from June 10th to June 15th, 2019. The conference will consist of one day of tutorials (June 10), followed by three days of main conference sessions (June 11-13), followed by two days of workshops (June 14-15).

\rule{0.48\textwidth}{0.4pt}
{\bf Megatron-LM}: With a broad scope, the conference addresses the challenges and opportunities in machine learning for practitioners and researchers. The conference program will feature sessions that cover different machine learning topics, from reinforcement learning to computational biology and from AI to deep learning. It is the world's premier forum for exploring and presenting advances in machine learning.
\\[1\baselineskip]
ICML 2019 will offer an impressive range of sessions, which will be of interest to practitioners, researchers, and students. Some of the topics we will cover include:
\\[1\baselineskip]
Reinforcement learning
\\[1\baselineskip]
Soft Computing
\\[1\baselineskip]
Computational biology
\\[1\baselineskip]
Convolutional networks
\\[1\baselineskip]
Deep learning
\\[1\baselineskip]
Statistical learning theory
\\[1\baselineskip]
Statistical learning machine
\\[1\baselineskip]
Natural language processing
\\[1\baselineskip]
Biology
\\[1\baselineskip]
Awards
\\[1\baselineskip]
Since its first edition, the conference has grown and evolved. Today, it is the premier forum for showcasing state-of-the-art research in machine learning. The conference is widely recognized as the preeminent venue for showcasing the latest research in machine learning, and a place where practitioners, researchers, and students can exchange ideas and insights.
\\[1\baselineskip]
The conference is a peer-reviewed venue that draws its prestige and influence from the excellence of the presentations and the quality of the invited speakers. With over 50\% of all talks selected based on their content by an expert committee, it is clear that the conference's high standards of research and quality keep a very high bar.
%\newline \newline
%For more information about the conference, please visit: ICML.
}}


\fbox{%
\parbox{0.48\textwidth}{%
\rule{0.48\textwidth}{0.4pt}
{\bf Context}: Rescue workers on Tuesday found a flight recorder from the Russian plane that crashed into the Black Sea over the weekend, the defense ministry said. All 92 people aboard the Russian military's Tu-154 plane are believed to have died Sunday morning when it crashed two minutes after taking off from the southern Russian city of Sochi.

\rule{0.48\textwidth}{0.4pt}
{\bf Megatron-LM}: A child holds up a portrait of a dead relative at a makeshift morgue in a building belonging to a local government in the town of Petropavlovsk-Kamchatsky. A Russian Defense Ministry Tu-154 plane with 92 people on board crashed into the Black Sea on Sunday morning shortly after taking off from the southern Russian city of Sochi. (Yuri Kochetkov/EPA)
\\[1\baselineskip]
According to the ministry, the plane was found at a depth of more than 1,000 feet and the recorder was damaged. "The defense ministry has decided to hand over the recorder to the Interstate Aviation Committee, which is the main investigative body in the air crash," the statement said.
\\[1\baselineskip]
The search for the plane's black boxes is being hampered by strong winds and rough seas. Investigators have launched a criminal probe into the crash, which officials said was likely caused by a technical fault.
\\[1\baselineskip]
Russian President Vladimir Putin declared Wednesday a day of mourning for the victims.

}}